\documentclass{article}

\usepackage{tabularx}
\usepackage{booktabs}

\title{Problem Statement and Goals\\\progname}

\author{\authname}

\date{}

\input{../Comments.text}
\input{../Common.text}

\begin{document}

\maketitle

\begin{table}[hp]
\caption{Revision History} \label{TblRevisionHistory}
\begin{tabularx}{\textwidth}{llX}
\toprule
\textbf{Date} & \textbf{Developer(s)} & \textbf{Change}\\
\midrule
Sep 27, 2026 & John Cho & Initial Problem Statement\\
Sep 27 2026 & Vikram Chandar & Goals, Stretch Goals and Environment\\
\bottomrule
\end{tabularx}
\end{table}

\newpage

\section{Problem Statement}

% \wss{You should check your problem statement with the
% \href{https://smiths.github.io/capTemplate/Checklists/ProbState-Checklist.pdf}
% {problem statement checklist}}. 

% \wss{You can change the section headings, as long as you include the required
% information.}

\subsection{Problem}

% \wss{What problem are you trying to solve?  Why is it important? Do not talk
% about how you will solve the problem, only what the problem is.  AI is rarely
% part of the problem.  The problem is ``what'' you want to do, not ``how'' you
% want to do it.  The problem statement should be written in a way that is
% understandable to a non-technical audience.}
Learning to fly a FPV (First person view) drone can be difficult and unsafe especially 
for beginner pilots. If left unaccounted for, they can risk damaging their drones, nearby 
property, or other equipment while learning basic flight controls and maneuvers.
\\ \\
Practing and flying in general through traditional means can also be limited by the need 
for large open spaces, weather conditions, and access to costly hardware. 
These barriers can make it difficult for anybody to fly a drone with little to no
consequences.

\subsection{Inputs and Outputs}

\textbf{Inputs:}
\begin{itemize}
    \item Joystick controller or similar interface (over RF or Serial)
\end{itemize}
\textbf{Outputs:}
\begin{itemize}
    \item Simulated feedback from the VR
    \item AI-integrated performance analysis at the end of a course/session
\end{itemize}

\subsection{Stakeholders}
\subsubsection{Primary Stakeholders}
\begin{itemize}
    \item \textbf{Drone Pilots:} FPVR will be designed with beginner to intermediate
    level pilots in mind, licensed or not, who have yet to familiarize themselves with flying drones
    in the field.
    \item \textbf{Drone Training Organizations:} In Canada, any drone above 250g in weight
    requires a pilot with a formal certificate to operate in a public space [1]. 
    Realistic drone flight simulation may complement these
    mandatory in-person flight exams facilitated by drone schools and provide pilots an easily accessible
    way to practice.
\end{itemize}

\subsubsection{Secondary Stakeholders}
\begin{itemize}
    \item \textbf{Drone Clubs and Teams:} Local technical clubs designing and building
    teleoperated FPV drones can benefit from FPVR and provide quick and useful feedback for improving
    the product. 
    \item \textbf{FPV Communities:} Anyone from a passive to hardcore FPV hobbyist
    within a drone community might find this technology useful both for casual and practical use. 
\end{itemize}

\subsection{Environment}


\subsubsection{Hardware Environment}

The simulator will run on a compatible VR headset capable of supporting real-time 3D applications. The headset will provide the display and motion tracking required for the user to interact with the simulated environment. VR controllers or compatible FPV flight controls will be used to control the simulated drone.

\subsubsection{Software Environment}

The simulator will run on the VR headset using a VR-compatible operating system and runtime environment. The software will require the necessary VR support to provide real-time graphics, motion tracking, and user input. The simulator will operate as a standalone VR application on supported hardware.


\section{Goals}

\begin{enumerate}
    \item \textbf{Provide a Realistic FPV Drone Flight Experience:} Create an immersive virtual reality FPV drone simulation that maintains a smooth and consistent frame rate during normal gameplay.

    \item \textbf{Support Beginner and Advanced Flight Training:} Enable new users to develop fundamental FPV flight skills while allowing experienced pilots to refine their performance through improvements in course completion times, flight consistency, and crash frequency.

    \item \textbf{Provide Realistic and Responsive Flight Controls:} Deliver intuitive controls with minimal input delay and realistic, consistent drone responses to throttle, pitch, roll, and yaw inputs.

    \item \textbf{Provide Varied Flight Environments and Courses:} Offer multiple environments and courses featuring different layouts, obstacles, and flight challenges to provide varied training scenarios.

    \item \textbf{Track Flight Performance:} Record basic flight performance data, including completion times, crashes, and course completion status, allowing users to review and monitor their progress.
\end{enumerate}

\section{Stretch Goals}

\begin{enumerate}
    \item \textbf{Provide Multiple Drone Configurations:} Offer multiple drone configurations with distinct flight characteristics, handling, and responsiveness, allowing users to experience different drone setups and adapt their flying techniques accordingly.

    \item \textbf{Support Multiplayer Flight:} Enable multiple users to operate drones within the same virtual environment, allowing them to practice and interact together in shared training scenarios.

    \item \textbf{Provide AI-Based Flight Analysis and Feedback:} Analyze controller inputs and flight performance to identify patterns such as inefficient inputs, inconsistent control, excessive corrections, and other areas for improvement, and provide personalized feedback to help users refine their flying skills.
\end{enumerate}

\section{Custom Documents (CDs)}

The extras chosen are the Design Thinking Report and the Usability Report. The Design Thinking Report will explain the intent behind our high level decisions and direction behind our design. The usability report will show the results of usability testing while also showing the response received from users. Our goals and methods regarding the testing will be documented as well.

\newpage{}

\section*{Appendix --- Reflection}

% \input{../Reflection.text}

\subsection*{John - Reflection}
\subsubsection*{What went well while writing this deliverable?}

Writing the problem statement itself was fairly simple since our
project was purposely selected to address a very specific use case
and audience. Our first TA meeting also helped to clarify certain sections of
the problem statement such inputs and outputs and how that relates to controller
and VR headset that we're planning on incorporating into the project.

\subsubsection*{What pain points did you experience during this deliverable, and how did you resolve them?}

The content of the report itself wasn't difficult to write, but coordinating who does what was a
much bigger headache than I expected. Although we had a solid grasp on how the project
will be laid out, we had no idea how to actually divide the work effectively. In the end,
we decided to manage all our tickets through a kanban and assign each of the tasks accordingly,
and raise any concerns/bugs through direct messages.

\subsubsection*{How did you and your team adjust the scope of your goals to ensure
    they are suitable for a Capstone project (not overly ambitious but also of
    appropriate complexity for a senior design project)?}

We mostly sorted all of that through in our TA meeting. There wasn't a lot that
we changed regarding the scope of the project, more so refined it to better suit
the user's needs. Originally, we were planning on designing this application purely 
for drone pilots to excerise and practice their skills, but soon realized we could
expand our target audience to include casual users who are more interested in the
gameplay.

\subsection*{Vikram - Reflection}
\subsubsection*{What went well while writing this deliverable?}

We were able to clearly define the project's goals, requirements, and technical direction. Discussing the available tools and technologies helped us identify which existing solutions could be used and where custom implementations would be necessary. We also found that dividing the work between team members made it easier to research different technical areas.

\subsubsection*{What pain points did you experience during this deliverable, and how did you resolve them?}

One of the main challenges was determining which technologies and features were realistic to implement within the time and resources available to the team. There were many possible tools and features that could improve the project, but including all of them would have increased the scope significantly. We addressed this by discussing the requirements as a team and prioritizing the core functionality needed for the drone simulation and VR training experience. 

\subsubsection*{How did you and your team adjust the scope of your goals to ensure
    they are suitable for a Capstone project (not overly ambitious but also of
    appropriate complexity for a senior design project)?}

The team narrowed the project around the core goal of creating a VR-based FPV drone training simulator with realistic flight controls and performance evaluation. Instead of attempting to build every feature from the beginning, we separated the project into core requirements and stretch goals. The core system will focus on the drone simulation, VR controls, training environments, performance tracking, and a basic flight evaluation system. Features such as multiple drone configurations, multiplayer, and more advanced AI-based analysis will be developed only if the core requirements are completed successfully.

\begin{thebibliography}{00}
    \bibitem{1}
        T. Canada, “Flying your drone safely and legally,” Transport Canada, Nov. 04, 2025. 
        https://tc.canada.ca/en/aviation/drone-safety/learn-rules-you-fly-your-drone/flying-your-drone-safely-legally
\end{thebibliography}
\end{document}